\chapter{Professional Communication ;  Written}

\section{Engineering Report Structure}
Every major deliverable (SOW, PDR report, CDR report, FDR report) follows a similar structure: \textbf{executive summary} (the most-read and least-written section; it should stand alone as a complete summary (write it last), \textbf{introduction} (context, scope, purpose), \textbf{body sections} organized by topic (not by team member), \textbf{conclusions and recommendations}, \textbf{references} (IEEE format), and \textbf{appendices} (detailed data, drawings, datasheets).

The executive summary deserves special attention. A busy client or reviewer may read only the executive summary. It must convey: what the project is, what was accomplished, what the key results are, and what the recommendations are) in one page or less. If the executive summary is unclear, the report has failed regardless of what follows.

\section{Technical Writing Best Practices}
\textbf{Active voice}: ``The team measured the output voltage'' (clear subject and action) not ``The output voltage was measured'' (who measured it?). Active voice is shorter, clearer, and more accountable.

\textbf{Specific language}: replace vague words with numbers and units. ``The temperature rose by 12.3\,\degC{} over 45 seconds'' not ``The temperature rose significantly.'' ``Significantly'' conveys no engineering information.

\textbf{Figures and tables}: every figure and table must be \textbf{numbered}, \textbf{captioned} (the caption should make the figure self-contained (a reader should understand it without reading the surrounding text), and \textbf{referenced in the text} (``as shown in Figure~3''). Never include a figure without discussing it. Never discuss a result without showing the data.

\textbf{Document organization}: organize by topic, not by team member. The reader does not care who wrote which section; they care about the logical flow of the engineering argument. Use consistent formatting, heading hierarchy, and terminology throughout.

\section{Citation and Reference Practices}
Use \textbf{IEEE citation format}. Cite all sources. Attribute all images, data, diagrams, and content not produced by the team. ``Hallucinated'' AI citations) references that do not actually exist; are an \textbf{academic integrity violation}. Every reference must be real and verified by clicking the link or locating the source.

\section{Document Version Control}
Engineering documents evolve throughout the project. Without version control, you end up with five different versions of the SDRD in five different folders and nobody knows which is current. Implement a simple version control system from Sprint~1:

\textbf{File naming}: use a consistent convention. Example:

\smallskip
\centerline{\texttt{TeamXX\_SDRD\_v2.1\_2026-03-15.docx}}
\smallskip

\noindent where XX is your team number, v2.1 is the version (major.minor), and the date is ISO format. The major version increments at each design review baseline (v1 = PDR baseline, v2 = CDR baseline, v3 = FDR baseline); the minor version increments for internal revisions.

\textbf{Change tracking}: use track changes (Word) or revision history (Google Docs) when editing team documents. Before submitting, review all tracked changes and accept or reject them as a team. The final submitted version should not contain visible track changes.

\textbf{Single source of truth}: designate one location (Teams folder, Google Drive folder) as the authoritative repository. All other copies are working drafts. When in doubt about which version is current, check the repository. The Documentation Lead is responsible for keeping the repository organized and current.

\section{Common Writing Errors in Capstone Reports}
These patterns appear in nearly every semester of capstone deliverables:

\textbf{Passive voice overuse}: ``The temperature was measured'' (who measured it?). ``The analysis was performed'' (by whom?). Active voice is clearer, shorter, and more accountable: ``The team measured the temperature'' or ``The structural analysis showed a safety factor of 2.3.'' Use active voice as the default; passive voice only when the actor is truly unknown or irrelevant.

\textbf{Imprecise language}: ``The system performed well'' (by what metric?). ``The results were close to the requirement'' (how close?). Replace subjective assessments with numbers: ``The system measured temperature with $\pm$0.7\,\degC{} accuracy, meeting the $\pm$1.0\,\degC{} requirement.''

\textbf{Unsupported claims}: ``This is the best approach'' (according to whom?). ``Industry standard practice is\ldots'' (cite the standard). Every factual claim must be supported by either your own data (test results, analysis) or an external source (textbook, standard, peer-reviewed paper, manufacturer datasheet).

\textbf{Missing transitions}: sections that jump from topic to topic without connecting paragraphs. The reader should understand how each section relates to the next. A one-sentence transition (``Having established the requirements, we now describe the concept generation process'') maintains narrative flow.

\textbf{Orphaned figures}: a figure appears in the document but is never referenced in the text. Every figure must be numbered, captioned, and discussed in the body text (``As shown in Figure~3, the output voltage increases linearly with temperature over the calibrated range'').

\textbf{Disorganized appendices}: critical data buried in appendices that are not referenced from the main text. Appendices supplement the main text; they do not replace it. Reference each appendix from the section that relies on it (``Complete test data is provided in Appendix~B'').

\textbf{Inconsistent terminology}: calling the same component ``the sensor module,'' ``the measurement unit,'' and ``the sensing element'' in different sections. Pick one term and use it throughout the document. Define it in the glossary if it is technical.

\section{The Report Revision Process}
Writing a professional engineering report is an iterative process, just like the design process itself. Plan for multiple revision cycles:

\textbf{First draft} (individual contributions): each team member writes their assigned sections. Focus on content completeness, not polish. Get the engineering substance on paper.

\textbf{Integration and peer review} (team activity): assemble all sections into a single document. Read the entire report as a team. Check for: consistent terminology, logical flow between sections, complete cross-references (figures, tables, sections), and gaps in content. Assign one person as the integration lead who reviews the entire document for consistency.

\textbf{Technical review}: ask your PA (or a teammate who did not write a particular section) to read critically. Can they understand the engineering argument without asking clarifying questions? If not, the section needs revision.

\textbf{Final polish}: proofread for grammar, spelling, formatting, and compliance with the Program Format Guide. Read the document aloud; errors that are invisible when reading silently become obvious when read aloud.

\textbf{Timeline}: the first draft should be complete 7--10 days before the submission deadline. Integration and peer review: 3--5 days before. Final polish: 1--2 days before. A report written in the final 48 hours will read like a report written in 48 hours.

\section{Report Template and Organization}
The Capstone Design Program Format Guide on Canvas specifies formatting requirements. Beyond the format, organize your report for readability:

\textbf{Section numbering}: use hierarchical numbering (1, 1.1, 1.1.1). Do not go deeper than three levels; if you need a fourth level, the section is too granular and should be restructured.

\textbf{Executive summary}: write this last, not first. It should stand alone as a complete summary (a reader who reads only the executive summary should understand: what the project is, what was accomplished, what the key results are, and what the recommendations are. One page maximum.

\textbf{Introduction}: set the context. What is the project? Who is the client? What problem does the system solve? This is not a copy of the executive summary; it provides the detailed background that the executive summary condenses.

\textbf{Body sections}: organize by topic, not by team member. The reader does not care who wrote which section; they care about the logical flow of the engineering argument. Within each section, follow the pattern: introduce the topic, present the methodology, show the results, and draw conclusions.

\textbf{Appendices}: use appendices for detailed data that supports the main text but would interrupt the narrative flow if included inline. Always reference appendices from the main text (``Complete test data is provided in Appendix~B'').

\section{Visual Communication in Reports}
Engineering reports communicate through visuals as much as text. Best practices:

\textbf{Figures}: use figures to show what words cannot efficiently describe) system architecture, test setups, data trends, CAD renderings, fabrication photos. Every figure must be numbered, captioned, and referenced in the text. The caption should make the figure self-contained.

\textbf{Tables}: use tables for structured data comparisons; requirements tables, BOM summaries, test result summaries, material properties. Every table must be numbered and captioned.

\textbf{Data visualization}: choose the right plot type. Line charts for trends over time or continuous variables. Bar charts for categorical comparisons. Scatter plots for correlation between two variables. Box plots for distribution comparisons. Never use 3D charts, pie charts with many slices, or gratuitous color when a simple plot would suffice.

\textbf{Consistent formatting}: use the same font, line weight, color scheme, and labeling convention throughout the document. Inconsistent formatting signals a document assembled hastily from disparate sources.



\section{The Ethics Assignment: Structure and Common Mistakes}
The Ethics Assignment (SDI Sprint~4) is a \textbf{1000--1500 word argumentative essay}; not an informative report, not a book summary, and not a list of ethical principles. You must:

\begin{enumerate}[nosep,leftmargin=1.5em]
\item \textbf{Take a controversial stance} on an ethical issue directly connected to your specific project's hardware, software, or system. ``Engineers should be ethical'' is not controversial. ``Autonomous irrigation systems should be prohibited in drought-restricted areas even when they reduce total water consumption'' is controversial and specific.
\item \textbf{Defend that stance with evidence}: cite data, case studies, regulations, or expert opinions that support your position. At least 3 credible sources.
\item \textbf{Address opposing viewpoints honestly}: present the strongest argument against your position and explain why your position is still preferable. This demonstrates critical thinking, not just advocacy.
\item \textbf{Write with a distinct voice}: the essay should sound like you, not like a generic AI-generated summary of ethical principles. Personal engagement with the specific ethical tension in your project is what distinguishes excellent essays.
\end{enumerate}

\textbf{The 7 most common mistakes} (avoid all of them):
\begin{enumerate}[nosep,leftmargin=1.5em]
\item Writing an informative essay instead of an argumentative one (no thesis, no position taken).
\item Taking a position so obvious that no reasonable person would disagree (``engineers should not build unsafe products'').
\item Discussing ethics in general instead of connecting to your specific project.
\item Failing to address counterarguments.
\item Exceeding the word limit (1500 words maximum; concision is a skill).
\item Using AI to generate the essay without genuine personal engagement.
\item Modifying the provided template formatting (use the template exactly as provided).
\end{enumerate}

\section{Figure and Table Formatting Standards}
Every figure and table in your reports must comply with these standards:

\textbf{Figures}: numbered sequentially (Figure~1, Figure~2, ...). Caption below the figure. Caption is a complete sentence describing what the figure shows and why it matters. All axes labeled with units. Legend if multiple data series. Referenced in the body text before the figure appears (``As shown in Figure~3, the output voltage increases linearly with temperature'').

\textbf{Tables}: numbered sequentially (Table~1, Table~2, ...). Caption above the table. Column headers with units. Consistent decimal places within each column. Referenced in the body text. Use horizontal rules (top, header, bottom) per the booktabs convention; avoid vertical rules and excessive gridlines.

\textbf{Equations}: numbered if referenced later. Variables defined immediately after the equation (``where $V$ is voltage in volts, $I$ is current in amperes, and $R$ is resistance in ohms''). Units stated explicitly.



\section{Collaborative Writing Workflows}
Capstone reports are written by teams, not individuals. This creates coordination challenges that must be managed:

\textbf{The outline-first approach}: before anyone writes prose, the team agrees on a complete outline: section titles, subsection titles, key content points for each section, and the responsible author. This prevents overlap (two people writing about the same topic) and gaps (no one writing about a critical topic).

\textbf{The integration meeting}: after first drafts are complete, hold a 2-hour team meeting where the integration lead reads the entire document aloud. The team listens for: inconsistent terminology (the sensor is called ``the thermocouple module'' in one section and ``the temperature sensor'' in another), gaps in logic (a result is referenced but never presented), redundant content (the same information appears in two sections), and tone inconsistencies (one section is formal, another is casual).

\textbf{The single-voice edit}: after integration, one team member (preferably the strongest writer) does a final ``voice pass'' through the entire document to make it read as if written by a single author. This pass fixes: inconsistent formatting, varying heading hierarchies, different citation styles, and tone shifts between sections.

\textbf{Deadline discipline}: the team writing schedule should work backward from the submission deadline. Example for a PDR report due Friday:
\begin{itemize}[nosep,leftmargin=1.5em]
\item Monday (5 days before): individual section drafts due to shared repository.
\item Tuesday: integration meeting; assemble sections, identify gaps and conflicts.
\item Wednesday: gap-filling and revision of flagged sections.
\item Thursday: single-voice edit, formatting check, PA review (if time permits).
\item Friday morning: final proofread, PDF generation, submission.
\end{itemize}

A team that starts writing on Thursday night for a Friday submission will produce a document that reads like it was written Thursday night.

\section{The Capstone Design Program Format Guide}
All deliverables should follow the Format Guide available on Canvas unless superseded by specific client requirements. The guide specifies: file format, heading styles, margin widths, citation format, and deliverable naming conventions. Comply with your PA and client preferences for any additional formatting requirements.

Key formatting rules that are frequently violated:
\begin{itemize}[nosep,leftmargin=1.5em]
\item \textbf{Page numbers}: every page numbered. Front matter uses roman numerals; body uses arabic numerals.
\item \textbf{Heading hierarchy}: consistent throughout (Title Case or Sentence case, not both).
\item \textbf{Figure/table numbering}: sequential within the document. Every figure and table referenced in the text.
\item \textbf{Citations}: IEEE format. Numbered in order of first appearance, not alphabetical.
\item \textbf{File naming}: use the convention specified by your PA. Typical format:

\smallskip
\centerline{\texttt{TeamXX\_[Deliverable]\_[Version]\_[Date].pdf}}
\end{itemize}

\section{The Writing Center}
The Writing Center at Colorado School of Mines (\url{https://writing.mines.edu}) provides free support for technical reports, essays, and presentations. Make an appointment at \url{https://mines.mywconline.com}. This is one of the most underutilized and most valuable campus resources. The Writing Center staff are trained to help engineers communicate more effectively; they understand that engineering writing has different conventions than humanities writing.



\section{The Final Proofread: A Systematic Approach}
The difference between a B+ report and an A report is often the final proofread. Use this systematic approach:

\textbf{Pass 1: Structure check}: read only the headings and topic sentences. Does the document have a logical flow? Does each section follow naturally from the previous one? Are there gaps (topics mentioned but never explained) or redundancies (the same content in two places)?

\textbf{Pass 2: Technical accuracy}: read every number, unit, equation, and figure caption. Does every number have appropriate significant figures and units? Does every equation define its variables? Does every figure caption describe what the figure shows? Do the numbers in the text match the numbers in the figures and tables?

\textbf{Pass 3: Citation check}: does every factual claim have a citation? Is every citation formatted correctly (IEEE numbered format)? Does every reference in the bibliography actually exist (no AI-hallucinated citations)? Is every figure and table cross-referenced in the text?

\textbf{Pass 4: Grammar and style}: read the document aloud (or use text-to-speech). Awkward sentences, run-on sentences, and missing words become obvious when heard. Check for: subject-verb agreement, consistent tense (past tense for completed work, present tense for ongoing truths), and parallel construction in lists.

\textbf{Pass 5: Format compliance}: check against the Capstone Design Program Format Guide. Page numbers present? Margins correct? Heading styles consistent? File named per the convention? All required sections present?

This five-pass approach takes 2--3 hours for a 30-page report. It is the highest-ROI activity you can perform before submission.



\section{Common Team Writing Pitfalls}
These patterns recur in capstone team reports every semester:

\textbf{The ``assembled, not integrated'' report}: each team member writes their section independently, the sections are concatenated into a single document, and nobody reads the result. The report has: three different names for the same component (``the sensor module'' in Chapter~2, ``the measurement unit'' in Chapter~4, ``the sensing subsystem'' in Chapter~6), two sections that explain the same concept differently, an introduction that promises content the body does not deliver, and formatting that changes between sections. Fix: the integration meeting and single-voice edit (described earlier) prevent this.

\textbf{The ``requirements rewrite'' report}: the team copies the SDRD into the report as Chapter~3, then rewrites the same requirements in different words throughout the design, analysis, and testing chapters. The result is a document where the reader encounters the same requirement 4--5 times in slightly different phrasing, creating confusion about which version is authoritative. Fix: present requirements once (in the SDRD chapter), then reference them by ID in subsequent chapters (``As specified in PR-004, the system shall maintain $\pm$1\,\degC{} accuracy'').

\textbf{The ``screenshot dump'' appendix}: the team includes 40 pages of unprocessed screenshots (oscilloscope captures without annotations, CAD models without dimensions, code output without interpretation. Appendix content must be organized, labeled, and referenced from the main text. Every figure in the appendix should have a caption explaining what it shows and why it matters. An appendix that nobody can navigate is an appendix that adds volume but not value.

\textbf{The ``written at midnight'' conclusion}: the final chapter is visibly rushed) shorter than preceding chapters, full of vague statements (``the project was successful''), and lacking the specificity of the rest of the document. The conclusion and recommendations chapter is the last thing the reviewer reads and strongly influences their overall impression. Write it with the same rigor as the analysis chapter.

\textbf{The missing ``so what''}: the report presents data but does not interpret it. ``Figure~7 shows the temperature response over 60 seconds'' tells the reader what the figure contains. ``Figure~7 shows that the system reaches the setpoint within 12 seconds with 4\% overshoot, meeting the PR-003 requirement of $\leq$15 seconds settling time and $\leq$10\% overshoot'' tells the reader what the data \emph{means}. Every data presentation must include engineering interpretation.



\section{IEEE Citation Format: Quick Reference}
Capstone reports use IEEE citation format. Key conventions:

\textbf{In-text citations}: numbered in square brackets in order of first appearance: ``The Nyquist theorem requires sampling at twice the signal frequency [1].'' Multiple citations: ``Several studies confirm this approach [2], [5], [7].'' Citation ranges: ``[3]--[6].''

\textbf{Reference list format} (common types):

\textit{Journal article}: [1] A.\,B.\ Author, ``Title of article,'' \textit{Journal Name}, vol.\ XX, no.\ YY, pp.\ ZZ--ZZ, Month Year.

\textit{Conference paper}: [2] A.\,B.\ Author, ``Title of paper,'' in \textit{Proc.\ Conference Name}, City, Country, Year, pp.\ ZZ--ZZ.

\textit{Book}: [3] A.\,B.\ Author, \textit{Title of Book}, Edition.\ City: Publisher, Year.

\textit{Datasheet}: [4] Manufacturer, ``Part Number Datasheet,'' Document Number, Year. [Online]. Available: URL.

\textit{Website}: [5] Organization, ``Page Title.'' URL (accessed Month Day, Year).

\textbf{Common citation mistakes}:
\begin{itemize}[nosep,leftmargin=1.5em]
\item Citing a URL without a title, author, or access date.
\item Citing Wikipedia as a primary source (cite the original sources that Wikipedia references).
\item Including references that are never cited in the text (or citing references not in the list).
\item AI-hallucinated citations that do not exist; \textbf{verify every reference by locating the actual source}.
\item Inconsistent formatting (some references have periods, others do not; some have volume numbers, others do not).
\end{itemize}



\section{Writing for Different Deliverable Types}
Each capstone deliverable has a distinct purpose and audience:

\textbf{The SOW}: primarily a \emph{contractual} document. Language is precise, commitments are explicit, and ambiguity is the enemy. Use SHALL/SHOULD for obligations. Define terms that could be interpreted differently by different readers. The SOW is not a design document; it defines \emph{what} will be delivered, not \emph{how}.

\textbf{The PDR/CDR/FDR reports}: primarily \emph{technical} documents. Language balances precision with readability. The narrative should flow logically: problem $\to$ approach $\to$ analysis $\to$ results $\to$ conclusions. Support every claim with evidence (data, calculations, references). The audience is technically sophisticated (PA, TAs, client engineers) and will challenge unsupported assertions.

\textbf{The Ethics and Broader Impacts essays}: primarily \emph{argumentative} documents. Language is persuasive but evidence-based. You are not reporting facts; you are defending a position. Use first person where appropriate (``I argue that...''), acknowledge counterarguments honestly, and demonstrate personal engagement with the ethical tension.

\textbf{The O\&M manual}: primarily an \emph{instructional} document. Language is simple, direct, and action-oriented. Use imperative mood (``Press the START button. Wait for the green LED. Verify the display shows HOME.'') with numbered steps. The audience may not be an engineer; assume no technical background unless the client specifies otherwise.

\textbf{Meeting minutes}: primarily a \emph{record} document. Language is factual and concise. Document decisions (not discussions), action items (with owners and deadlines), and attendance. Meeting minutes are legal-weight documents in professional practice; in capstone, they are the evidence base for PIPs and grade disputes.


\section{Practice Questions}
\begin{enumerate}[leftmargin=1.5em]
\item Rewrite this sentence in active voice with specific data: ``It was found that the system performed adequately under the tested conditions.''
\item Write a one-paragraph executive summary for a hypothetical CDR report for your project.
\item A figure in your PDR report shows test data but has no caption, no axis labels, and no reference in the text. List everything that needs to be fixed to meet professional standards.
\end{enumerate}


